% steps/step13.tex
\section[گام سیزدهم: بازآموزی و بروزرسانی]{\faUpload\ \ گام سیزدهم: بازآموزی و بروزرسانی}
\begin{stepbox}

\field{معیارها و زمان نیاز به بازآموزی مدل}{
    نیاز به بازآموزی مدل زمانی قطعی می‌شود که شاخص‌های مانیتورینگ نشان‌دهنده افت کارایی مداوم باشند، خطاهای سیستماتیک در گزارش‌های پزشکان ثبت شود، یا اسکنرهای جدیدی در بیمارستان راه‌اندازی شوند که خروجی متفاوتی دارند. در این شرایط، چرخه‌ی آموزش مجدد مدل فعال می‌گردد.
}

\field{به‌روزرسانی داده‌ها و یادگیری فعال (\lr{Active Learning})}{
    تصاویری که مدل در دنیای واقعی در پردازش آن‌ها دچار خطا شده (مانند تشخیص اشتباه کیست‌ها به جای تومور) پس از اصلاح دستی توسط رادیولوژیست، به‌عنوان داده‌های ارزشمند به پایگاه داده آموزشی اضافه می‌شوند. با تمرکز بر روی نمونه‌های چالش‌برانگیز و سخت (\lr{Hard Negative Mining})، مدل جدید با کارایی بالاتر و با استفاده از منابع پردازشی کمتر بازآموزی می‌شود.
}

\field{مدیریت نسخه‌ها و عقب‌گرد مدل (\lr{Model Versioning \& Rollback})}{
    نسخه‌های جدید مدل نباید مستقیماً جایگزین نسخه در حال خدمت شوند. با استفاده از سیستم‌های مدیریت مدل (مانند \lr{MLflow})، هر نسخه شناسنامه و تاریخچه مشخصی دارد. پس از آموزش، نسخه جدید به‌صورت تدریجی مستقر شده و در صورت بروز هرگونه رفتار پیش‌بینی‌نشده بالینی، سیستم بلافاصله به نسخه پایدار قبلی عقب‌گرد (\lr{Rollback}) داده می‌شود.
}

\end{stepbox}
